\newpage
\section{Appendix}

\subsection{Extending tests for convergence of number series}
% Cauchy Condensation Test
\subsubsection{Talking from Cauchy Condensation Test}
\fbox{\begin{minipage}{0.95\textwidth}
    \begin{thm}
        (Caucy Condensation Test) \textcolor{skycyan}{Tech toolkit} \\
        $\{a_n\}_{n}^{\infty} \searrow^{+}$ (i.e. $\cdots \geq a_n \geq a_{n-1} \geq \cdots \geq a_1 \geq 0$), then $\sum_{n=1}^{\infty} a_n$ and $\sum_{k=1}^{\infty} 2^k a_{2^k}$ have same convergence.
    \end{thm}
\end{minipage}}
\newline\newline
Connections betwwen \underline{Condensation test} and \underline{Integral test}: \\
\begin{equation*}
    \begin{split}
        \sum_{n=1}^{\infty} a_n < \infty &\iff \int_{1}^{\infty} f(x) dx = \infty \ \text{substitute with: } x=2^y \\
        &\iff \int_{1}^{\infty} 2^y f(2^y) dy < \infty \iff \sum_{k=1}^{\infty} 2^k a_{2^k} < \infty
    \end{split}
\end{equation*}
\fbox{\begin{minipage}{0.95\textwidth}
    \begin{thm}
        (Maclaurin-Cauchy Integral Test) \textcolor{skycyan}{Tech toolkit} \\
        $f \searrow^{+}$ on $[1, +\infty)$, then $\sum_{n=1}^{\infty} f(n)$ and $\int_{1}^{\infty} f(x) dx$ have same convergence.
    \end{thm}
\end{minipage}}
\newline\newline
Here are some examples we could use Condensation test to solve: \\
\fbox{\begin{minipage}{0.95\textwidth}
    \begin{enumerate}
        \item $\sum_{n=1}^{\infty} \frac{1}{n^p}$, $\sum_{n=2}^{\infty} \frac{1}{n\log^p n}$, $\sum_{n=2}^{\infty} \frac{1}{(\log n)^p}$ (combine Cauchy root test), $\sum_{n=2}^{\infty} \frac{\log n}{n^2}$
        \item $\sum_{n=16}^{\infty} \frac{1}{n(\log n)(\log\log n)(\log\log\log n)^p}$
        \item More general, let's consider $f(x) := n^{-a}(\log n)^{-b}(\log\log n)^{-c}$, $\sum f(n) < \infty, \ a>1$, when $a=1$ condensation transformation gives $\sum n^{-b}(\log n)^{-c}$ (seems like the logarithms "shift to the left") so when $a=1$, it converges as $b>1$. When $b=1$, the value of $c$ enters. \\
        This result readily generalizes: the condensation test, applied repeatedly, can be used to show that for $k=1,2,3,\cdots$, the generalized Bertrand series $\sum _{n\geq N}{\frac {1}{n\cdot \log n\cdot \log \log n\cdots \log ^{\circ (k-1)}n\cdot (\log ^{\circ k}n)^{\alpha }}}\quad \quad (N=\lfloor \exp ^{\circ k}(0)\rfloor +1)$ converges for $\alpha>1$ and diverges for $0 \leq \alpha \leq 1$. \\
        from Wikipedia: \href{https://en.wikipedia.org/wiki/Cauchy_condensation_test}{Cauchy Condensation Test}
    \end{enumerate}
\end{minipage}}
\newline\newline
Besides, the Cauchy condensation test is a partial case of following classical result deu to Schlomilch: \\
\fbox{\begin{minipage}{0.95\textwidth}
    \begin{thm}
        (Schlomilch) \\
        $\{a_n\}_{n}^{\infty} \searrow^{+}$, $\{u_n\}_{0}^{\infty} : u_0<u_1<\cdots\ \text{and}\ \frac{\Delta u_n}{\Delta u_{n-1}}\leq C$, then $\sum_{n=1}^{\infty} a_n$ and $\sum_{n=1}^{\infty} \Delta u_n a_{u_n} = \sum_{n=1}^{\infty} (u_{n+1}-u_n)a_{u_k}$ have same convergence. 
    \end{thm}
\end{minipage}}
\begin{nt*}
    1. You may note it that let $u_n = 2^n$, it's Cauchy condensation test. Moreover, you could also use $u_n = p^n,\ p>1$.
\end{nt*}

% Review the History of the tests for convergence of number series - their monotonicity assumption
\subsubsection{History of tests with monotonicity assumption for sequence}
Let's recall some convergence tests with monotonicity assumption for sequence: \\
\fbox{\begin{minipage}{0.95\textwidth}
    \begin{thm}
        (Ermakov) \\
        Given $f \searrow_{\geq 0} \in C(1,\infty)$, if $\lim_{t\to\infty} \frac{e^tf(e^t)}{f(t)} \leq q < 1$, then $\sum_{n=1}^{\infty} f(n) < \infty$; if $\lim_{t\to\infty} \frac{e^tf(e^t)}{f(t)} \geq 1$, then $\sum_{n=1}^{\infty} f(n)$ diverges.
    \end{thm}
    \begin{nt*}
        1. You could replace $e^t$ with an arbitrary $\varphi(t) \nearrow_{\geq 0}$ to get a family of tests. 
    \end{nt*}
\end{minipage}}
\newline\newline
\fbox{\begin{minipage}{0.95\textwidth}
    \begin{thm}
        (Abel-Olivier) \\
        Given $\{a_n\}_{n}^{\infty} : a_n \searrow_{\geq 0} \to 0$ (monotone null), if $\sum_{n=1}^{\infty} a_n < \infty$, then $na_n = o(1),\ n\to\infty$.
    \end{thm}
    \begin{nt*}
        1. Proof idea: Cauchy convergence criterion. \\
        2. The inverse is not ture: $a_n = \frac{1}{n\log n}$. 
    \end{nt*}
\end{minipage}}
\newline\newline
\fbox{\begin{minipage}{0.95\textwidth}
    \begin{thm}
        (Sapogov) \\
        Given $\{a_n\}_{n}^{\infty} : a_n \searrow_{\geq 0}$, then $\lim_{n\to\infty} a_n = 0$ i.f.f. $\sum_{n=1}^{\infty} (1-\frac{a_n}{a_{n+1}})$ as well as $\sum_{n=1}^{\infty} (\frac{a_{n+1}}{a_n}-1)$ diverges.
    \end{thm}
    \begin{nt*}
        1. Proof idea: $\sum_{n=1}^{\infty} b_n$ and $\sum_{n=1}^{\infty} \log(1-b_n)$ has same convergence. 
        2. It's a simple test for $n$-th term of a series apporoach to $0$. 
        3. Other Description: \\
        Given $\{a_n\}_{n}^{\infty} : a_n \nearrow_{\geq 0}$, if $\sum_{n=1}^{\infty} (1-\frac{a_n}{a_{n+1}})$ as well as $\sum_{n=1}^{\infty} (\frac{a_{n+1}}{a_n}-1)$ converges, then $a_n$ is bounded or diverges otherwise. 
    \end{nt*}
\end{minipage}}
\newline\newline
\fbox{\begin{minipage}{0.95\textwidth}
    \begin{thm}
        (Deddkind-du Bois Reymond) \\
        Def: $\{a_n\}_{n}^{\infty} \in \text{BV(bounded variation)} : \sum_{n=1}^{\infty} |a_{n+1}-a_n| < \infty$. \\
        1. Given $\{a_n\}_{n}^{\infty} \in \text{BV}$ also $a_n = o(1)$ and $\sum_{k=1}^{n}b_k < C,\ \forall n\in\mathbb{N}$, then $\sum_{n=1}^{\infty} a_nb_n < \infty$. \\
        2. Given $\{a_n\}_{n}^{\infty} \in \text{BV}$ and $\sum_{n=1}^{\infty}b_n < \infty$, then $\sum_{n=1}^{\infty} a_nb_n < \infty$. 
    \end{thm}
    \begin{cor}
        1. (Dirichlet) \\
        Given $a_n$ is monotone null and $\sum_{k=1}^{n}b_k < C,\ \forall n\in\mathbb{N}$, then $\sum_{n=1}^{\infty} a_nb_n < \infty$. \\
        Corillary. (Leibniz) \\
        Given $a_n$ is monotone null, then $\sum_{n=1}^{\infty} (-1)^{n}a_n < \infty$. \\
        2. (Abel) \\
        Given $a_n$ is bounded monotone and $\sum_{n=1}^{\infty} b_n < \infty$, then $\sum_{n=1}^{\infty} a_nb_n < \infty$. 
    \end{cor}
\end{minipage}}
\newline\newline
Before our formal work, let's introduce some basic more general monotonicity definitions: \\
\fbox{\begin{minipage}{0.95\textwidth}
    \begin{df}
        We call a non-negative null (that is, tending to zero at infinity) sequence ${a_k}$ weak monotone, written \textbf{WMS}, if for some positive absolute constant $C$ it satisfies $a_k \leq C a_n \ \forall k \in [n, 2n]$.
    \end{df}
\end{minipage}}
\newline\newline
To introduce a counterpart for functions, we will assume all functions to be defined on $(0, \infty)$, locally of bounded variation, and vanishing at infinity. \\
\fbox{\begin{minipage}{0.95\textwidth}
    \begin{df}
        We say that a non-negative function $f$ defined on $(0, \infty)$, is weak monotone, written \textbf{WM}, if $f(t) \leq Cf(x) \ \forall t \in [x, 2x]$.
    \end{df}
\end{minipage}}
\newline\newline
Denote $\Delta a_k = a_{k+1} - a_k$. \\
\fbox{\begin{minipage}{0.95\textwidth}
    \begin{df}
        A positive null sequence ${a_k}$ is called general monotone, written \textbf{GMS}, if it satisfies $\sum_{k=n}^{2n} |\Delta a_k| \leq C a_n$, for any $n$ and some absolute constant $C$.
    \end{df}
\end{minipage}}
\newline\newline
\fbox{\begin{minipage}{0.95\textwidth}
    \begin{df}
        We say that a non-negative function $f$ is general monotone, \textbf{GM}, if $\forall x \in (0, \infty)$, $\int_{x}^{2x} |df(t)| \leq Cf(x)$.
    \end{df}
\end{minipage}}
\newline\newline
\fbox{\begin{minipage}{0.95\textwidth}
    \begin{nt*}
        1. The class of quasi-monotone sequences, that is, ${a_k}$ such that there exists $\tau > 0$ so that $k^{-\tau} a_k \downarrow$, is a proper subclass of $GMS$. \\
        2. One of the simple basic properties of $GMS$ is $GMS \subsetneq WMS$. \\
        3. We just remark that if $f(\cdot)$ is a $GM$ function, then for $a_k = f(k)$ there holds ${a_k} \in GMS$. \\
        4. Also $GM \subsetneq WM$.
    \end{nt*}
\end{minipage}}

% Extension of the "monotone" tests
\subsubsection{Extension of the "monotone" tests}
Now, let's extend those aforementioned "monotone" tests for sequences of numbers and functions: \\
\fbox{\begin{minipage}{0.95\textwidth}
    \begin{thm}
        (Maclaurin-Cauchy $\uparrow$) \\
        Let $f$ be a WM function. Then the series $\sum_{k=1}^\infty f(k)$ and integral $\int_1^\infty f(t) dt$ converge or diverge simultaneously.
    \end{thm}
\end{minipage}}
\newline\newline
\fbox{\begin{minipage}{0.95\textwidth}
    \begin{thm}
        (Ermakov $\uparrow$) \\
        Let $f$ be a WM function and let $\varphi(t)$ be a monotone increasing, positive function having a continuous derivative and satisfying $\varphi(t) > t$ for all $t$ large enough. \\
        If for $t$ large enough $\frac{f(\varphi(t))\varphi'(t)}{f(t)} \leq q < 1$, then series $\sum_{n=1}^{\infty} f(n)$ converges, while if $\frac{f(\varphi(t))\varphi'(t)}{f(t)} \geq 1$, then series $\sum_{n=1}^{\infty} f(n)$ diverges.
    \end{thm}
\end{minipage}}
\newline\newline
\fbox{\begin{minipage}{0.95\textwidth}
    \begin{thm}
        (Schlomilch $\uparrow$) \\
        Let ${u_k}$ be an increasing sequence of positive numbers such that $u_{k+1} = O(u_k)$ and $u_k \to \infty$. \\
        Let $f$ be a WM function. Then both series $\sum_{k=1}^{\infty} f(u_k) \Delta u_k$ and $\sum_{k=1}^{\infty} f(u_{k+1}) \Delta u_k$ converge (or diverge) with $\int_1^{\infty} f(t) dt$. \\
        Let ${a_k}$ be a WMS. Then $\sum_{n=1}^{\infty} a_n$ converges if and only if the series $\sum_{k=1}^{\infty} \Delta u_k a_{u_k} = \sum_{k=1}^{\infty} (u_{k+1} - u_k) a_{u_k}$ converges.
    \end{thm}
\end{minipage}}
\newline\newline
\fbox{\begin{minipage}{0.95\textwidth}
    \begin{thm}
        (Abel-Olivier $\uparrow$) \\
        Let ${a_k}$ be a WMS. If series $\sum_{n=1}^{\infty} a_n$ converges, then $k a_k$ is a null sequence.
    \end{thm}
\end{minipage}}
\newline\newline
Analyzing one Dvoretzky’s result and its proof, we obtain the following statement: \\
\newline\newline
\fbox{\begin{minipage}{0.95\textwidth}
    \begin{thm}
        If ${a_k}, {b_k} \in WMS$, and the series are convergent and divergent respectively, then $\forall M > 1$ there $\exists$ infinitely many $R_j$,
        $R_j \to \infty$, s.t. $\forall k$ with $R_j \leq k \leq M R_j$, we have $a_k < b_k$.
    \end{thm}
\end{minipage}}
\newline\newline
Our next result is "dual" to the above Schlömlich-type extensions. \\
\fbox{\begin{minipage}{0.95\textwidth}
    \begin{df}
        An increasing sequence ${u_k}$ is called \textbf{lacunary} if $u_{k+1}/u_k \geq q > 1$.
    \end{df}
\end{minipage}}
\newline\newline
A more general class of sequences is the one in which each sequence can be split into finitely-many lacunary sequences, which we will write ${u_k} \in \Lambda$. (This is true i.f.f. $\sum_{j=1}^k u_j \leq C u_k$.) \\
In different terms, ${u_k} \in \Lambda$ is true i.f.f. there $\exists r \in \mathbb{N}$ s.t. $\frac{u_{k+r}}{u_k} \geq q > 1, \ k \in \mathbb{N}$. \\
\fbox{\begin{minipage}{0.95\textwidth}
    \begin{df}
        $\Delta a_{u_k} := a_{u_k} - a_{u_{k+1}}$
    \end{df}
\end{minipage}}
\newline\newline
\fbox{\begin{minipage}{0.95\textwidth}
    \begin{thm}
        Let ${a_k}$ be a non-negative WMS, and let a sequence ${u_k}$ be s.t. ${u_k} \in \Lambda$ and $u_{k+1} = O(u_k)$. Then $\sum_{k=1}^{\infty} a_k$,
        $\sum_{k=1}^{\infty} u_k |\Delta a_{u_k}|$, and $\sum_{k=1}^{\infty} u_k a_{u_k}$ converge or diverge simultaneously.
    \end{thm}
\end{minipage}}
\newline\newline
It's also possible to get equiconvergence results for an important case $u_n = n$, where the lacunarity is not suitable. (we need to proceed a smaller class than WMS)\\
\fbox{\begin{minipage}{0.95\textwidth}
    \begin{thm}
        Let ${a_k}$ be a GMS. Then series $\sum_{n=1}^{\infty} a_n$ and $\sum_k k| \Delta a_k|$ converge or diverge simultaneously.
    \end{thm}
    \begin{thm}
        Let $f$ be a $GM$ function. Then the integrals $\int_1^\infty f(t) dt \text{ and } \int_1^\infty t|df(t)| \text{ converge or diverge simultaneously}$.
    \end{thm}
\end{minipage}}
\begin{nt*}
    Negative type results: \\
    1. Sapogov type test cannot be true if ${b_k}$ (as well as ${1/b_k}$) is WMS. (A counterexample can be constructed against generalization of Abel's test.) \\
    2. As for extending the Leibniz test, it cannot hold without additional assumption of the boundedness of variation: just take $a_k = 1/\ln k$ everywhere except $n = 2^k$ where $a_n = 2/\ln n$. \\
    3. Is there exists a wider class that those tests could hold? (see below problem)
\end{nt*}
\begin{prb*}
    WMS is, in a sense, the widest class for which such tests are still valid.
\end{prb*}
We give some examples to show you the problem is seemly true. 
\begin{ex*}
    \begin{enumerate}
        \item An immediate natural extension of $WMS$ is the class defined by $a_k \leq C \sum_{n=k/2}^k \frac{a_n}{n}$, or a bit more general $a_k \leq C \sum_{n=[k/c]}^{[ck]} \frac{a_n}{n},\ c > 1$. \\
        The principle difference between $WMS$ and these classes is that the latter two allow certain amount of zero members, unlike $WMS$ that forbid even a single zero, i.e., $a_{n_0} = 0$ implies $a_n = 0$ for $n \geq n_0$. \\
        Putting zeros on certain positions, say $k = 2^n$, we easily construct a counterexample to show that the Cauchy condensation test cannot be valid for these classes nor its extensions. \\
        In a similar way, just letting a function $f$ to take non-zero values only close to integer points, one sees that the Maclaurin-Cauchy integral test may fail as well. 
        \item In conclusion, note that $WMS$ is a subclass of the broadly used $\Delta_2$-class: for which the doubling condition $a_{2k} \leq C a_k$ holds $\forall k \in \mathbb{N}$. \\
        We observe that assuming the doubling condition by no means can guarantee the above tests to be extended: \\
        To illustrate this, the easiest way is to consider Cauchy's condensation test. Take ${a_k}$ such that $a_k = \begin{cases} 2^{-k}, & k \neq 2^n; \\ n^{-2}, & k = 2^n. \end{cases}$ This sequence is doubling but not $WMS$. Obviously, $\sum_{n=1}^\infty 2^n a_{2^n} = \sum_{n=1}^\infty \frac{2^n}{n^2}$ diverges. 
    \end{enumerate}
\end{ex*}

% Overview
\subsubsection{Summary}
The main results may be summarized as the following statement: \\
\fbox{\begin{minipage}{1.00\textwidth}
    \begin{thm}
        Let $f(\cdot) \in WM$. Then the following series and integrals converge or diverge simultaneously:
        $$\int_{1}^{\infty} f(t) dt, \quad \sum_{k} f(k);$$
        $$\sum_{k} (u_{k+1} - u_k)f(u_k), \quad \text{provided} \quad u_k \uparrow, \quad u_{k+1} = O(u_k);$$
        $$\sum_{k} u_k f(u_k), \quad \sum_{k} u_k |f(u_{k+1}) - f(u_k)|, \quad \text{with} \quad u_k \text{ lacunary}, \quad u_{k+1} = O(u_k);$$
        $$\sum_{k} k |f(k + 1) - f(k)|, \quad \int_{1}^{\infty} t |df(t)|, \quad \text{provided} \quad f(\cdot) \text{ is } GM.$$
    \end{thm}
\end{minipage}}
\begin{ex*}
    Besov space: \\
    For $k > r$, $\theta, r > 0$ and $p \geq 1$ $\|f\|_{B^r_{p,\theta}} = \left( \int_0^1 t^{-r\theta-1} \omega_k^\theta (f;t)_p \frac{dt}{t} \right)^{1/\theta}$. \\
    Equivalent form: $\|f\|_{B^r_{p,\theta}} = \left( \int_1^\infty t^{r\theta+1} \omega_k^\theta (f;1/t)_p \frac{dt}{t} \right)^{1/\theta}$. \\
    Since $\omega_k(f; t)p \uparrow$, the function under the integral sign is $WM$. Applying generalized Maclaurin-Cauchy and Cauchy condensation tests, we obtain convenient equivalent form: $\left( \sum_{\nu=0}^{\infty} \left[ 2^{r\nu} \omega_k(f; 2^{-\nu})_p \right]^{\theta} \right)^{1/\theta}$.
\end{ex*}

\newpage
\subsection{Review number series through exercises}
\subsubsection{Basic simple problems}
First we talk about some basic concepts of number series, some results could be used in the following exercises and most of them is the basics of complex and difficult problems. \\
\begin{prb*}
    1. Given $\sum a_n$ and $\sum b_n$ converge at same time, consider convergence of following series: 
    $$\sum (a_n\pm b_n),\quad \sum a_nb_n,\quad \sum \frac{a_n}{b_n}(b_n\neq 0),\quad \sum \max\{a_n,b_n\},\quad \sum \min\{a_n,b_n\}$$
    What happens if the series is all positive? \\
    2. If $\sum (a_n+a_{n+1})$ con., could we get $\sum a_n$ con.? What happens if $\sum a_n$ is all positive?
\end{prb*}
\fbox{\begin{minipage}{0.95\textwidth}
    \begin{sol*}
        1. Easy, just check for a warm-up every time. \\
        Here are some ideas or directions provided: 1). consider positive series and alternating series 2). consider the oder of terms of series (ex: $\sum \frac{1}{n^p}$) 3). in the process of constructing a series, some divergent series terms can be inserted into the convergent series, or vice versa. \\
        2. 1). $\sum (-1)^n$ (idea: $\sum a_{2k}\to\infty$, $\sum a_{2k+1}\to-\infty$, $a_{2k}\sim a_{2k+1}$) 2). $\sum a_n < \sum a_n+a_{n+1}$
    \end{sol*}
\end{minipage}}

\begin{prb*}
    If $\forall p\in\mathbb{Z}^+$, $\lim_{n\to\infty} (a_{n+1}+a_{n+2}+\cdots+a_{n+p}) = 0$, then what about $\sum a_n$?
\end{prb*}
\fbox{\begin{minipage}{0.95\textwidth}
    \begin{sol*}
        Core: the $N_p$ is different from $N_\varepsilon$ of Cauchy's convergence criterion. \\
        Example: $\sum \frac{1}{n}$ div. $\sum \frac{1}{n^2}$ con.
    \end{sol*}
\end{minipage}}

\begin{prb*}
    Given $a_n \leq c_n \leq b_n$, $\sum a_n$ and $\sum b_n$ both con. or div. $\nRightarrow$ $\sum c_n$ con. or div.
\end{prb*}
\fbox{\begin{minipage}{0.95\textwidth}
    \begin{sol*}
        Core: the condition of the squeeze theorem of a sequence is the order relationship of the terms of the sequence, but the terms of a series are different from the terms of a partial sum sequence. \\
        Example: 1. $\sum \frac{(-1)^n}{n}$ con. and $\sum \frac{(-1)^{n+1}}{n}$ con. $\nRightarrow$ $\sum \frac{1}{n}$ 2. $\sum \frac{1}{n}$ div. and $\sum -\frac{1}{n}$ div. $\nRightarrow$ $\sum 0$ con.
    \end{sol*}
\end{minipage}}

\begin{prb*}
    1. Given $\sum a_n \to S$ con., if $\forall n$ exchange position of $a_{2n-1}$ and $a_{2n}$, prove the new series con. \\
    2. What happens if we exchange $a_{pn-p+1}, \cdots, a_{pn}$?
\end{prb*}
\fbox{\begin{minipage}{0.95\textwidth}
    \begin{sol*}
        1. Note new series as $T_n$, then $|T_n-S|=|T_n-S_n+S_n-S|\leq|T_n-S_n|+|S_n-S|\leq|a_{n}|+|a_{n_1}|+o(1) = o(1)$. \\
        2. Prove is similiar (almost same). Note: actually the inequality is loose.
    \end{sol*}
\end{minipage}}

\subsubsection{Positive series}
\begin{prb*}
    1. Given positive-term $\sum \frac{1}{a_n}$ con., then $\sum \frac{n}{a_1+\cdots+a_n}$ con. \\
    2. (Carleman) Given positive-term $\sum a_n$ con., then $\sum \sqrt[n]{a_1\cdots a_n} \leq e\sum a_n$.
\end{prb*}
\fbox{\begin{minipage}{0.95\textwidth}
    \begin{sol*}
        1. \textcolor{skycyan}{/Tech idea/: Monotonicity} \\
        1) Suppose $\{a_n\} \uparrow$, then $a_1+\cdots+a_{2n} \geq na_n \implies \frac{2n-1}{a_1+\cdots+a_{2n-1}}+\frac{2n}{a_1+\cdots+a_{2n}}\leq \frac{2n-1}{na_n}+\frac{2n}{na_n} < \frac{4}{a_n}$ \\
        2) Rearrange $\{a_n\}$ to $\{b_n\} \uparrow$, $\sum \frac{n}{a_1+\cdots+a_n} \leq \sum \frac{n}{b_1+\cdots+b_n}$ con. \\
        2. \begin{align*}
        \sum_{n=1}^{N} \sqrt[n]{a_1 a_2 \cdots a_n} &= \sum_{n=1}^{N} \frac{\sqrt[n]{(a_1)(2a_2)\cdots(na_n)}}{\sqrt[n]{n!}} \leqslant \sum_{n=1}^{N} \frac{a_1 + 2a_2 + \cdots + na_n}{n \sqrt[n]{n!}} \\
        &\leqslant e \sum_{n=1}^{N} \frac{a_1 + 2a_2 + \cdots + na_n}{n(n+1)} = e \sum_{n=1}^{N} \frac{1}{n(n+1)} \sum_{k=1}^{n} k a_k \\
        &= e \sum_{k=1}^{N} k a_k \sum_{n=k}^{N} \frac{1}{n(n+1)} = e \sum_{k=1}^{N} k a_k \left( \frac{1}{k} - \frac{1}{N+1} \right) \\
        &\leqslant e \sum_{k=1}^{N} a_k
        \end{align*}
    \end{sol*}
\end{minipage}}

\begin{prb*}
    Given 2 +-term $\sum a_n$ and $\sum b_n$, partial sum $S_n$ and $S'_n$, remainder $R_n$ and $R'_n$, satisfies: $a_n \sim b_n$, prove: \\
    1. If $\sum a_n$ and $\sum b_n$ con., then $R_n \sim R'_n$; 
    2. If $\sum a_n$ and $\sum b_n$ div., then $S_n \sim S'_n$.
\end{prb*}
\fbox{\begin{minipage}{0.95\textwidth}
    \begin{sol*}
        1. $\frac{R_n}{R'_n} = \frac{\sum_{n=N}^{M} a_n + \sum_{n=M}^{\infty} a_n}{\sum_{n=N}^{M} b_n + \sum_{n=M}^{\infty} b_n} (M>N, N\to\infty)$: $\sum_{n=N}^{M} a_n \sim \sum_{n=N}^{M} b_n,\ \sum_{n=M}^{\infty} a_n, \sum_{n=M}^{\infty} b_n = o(1)$, using Cauchy's convergence Crite: we could choose $M>N$, s.t. $\sum_{n=M}^{\infty} a_n = \sum_{n=N}^{M} a_n (1+o(1)), \sum_{n=M}^{\infty} b_n = \sum_{n=N}^{M} b_n (1+o(1))$. \\
        2. (Abel type THeorem - 1)
    \end{sol*}
\end{minipage}}

\begin{prb*}
    Calculation: 
    \begin{enumerate}
        \item $\sum \frac{1}{2^{\sqrt{n}}}$: $\frac{1}{3^{\sqrt{n}}} < \frac{1}{3^{\log n}} = \frac{1}{n^{\log 3}},\ n\to\infty$ con.
        \item $\sum \frac{1}{n^{1+\frac{1}{n}}}$: $\frac{1}{n^{1+\frac{1}{n}}} \sim \frac{1}{n}\ (n^{\frac{1}{n}}),\ n\to\infty$ div.
        \item $\sum [\frac{1}{n}-\log(1+\frac{1}{n})]$: $\log(1+\frac{1}{n}) = \frac{1}{n}+O(\frac{1}{n^2})$ con.
        \item $\sum \frac{n^{n-1}}{(2n^2+n+1)^{\frac{n-1}{2}}}$: $\frac{n^{n-1}}{(2n^2+n+1)^{\frac{n-1}{2}}} = (\frac{n^2}{2n^2+n+1})^{\frac{n-1}{2}} = (\frac{1}{2})^{\frac{n-1}{2}}(1+o(1))$ con.
        \item \checkmark $\sum \frac{n^{\log n}}{(\log n)^n}$: $\frac{n^{\log n}}{(\log n)^n}<\frac{1}{n^2}\ (n>13) \iff (\log n)^2+2\log n-n\log(\log n)<0$ use function and derivative. con. 
        \item $\sum \frac{e^{-(1+\frac{1}{2}+\cdots+\frac{1}{n})}}{n^p}$: $\sum_{k=1}^{n} \frac{1}{k} = \log n + \gamma + O(\frac{1}{n})$ then $p>0$ con., $p\leq 0$ div.
        \item \checkmark $\sum (a^{\frac{1}{n}} -\frac{b^{\frac{1}{n}}+c^{\frac{1}{n}}}{2})\ (a,b,c>0)$: $a^{\frac{1}{n}} -\frac{b^{\frac{1}{n}}+c^{\frac{1}{n}}}{2} = e^{\frac{\log a}{n}} - \frac{1}{2}(e^{\frac{\log b}{n}}+e^{\frac{\log c}{n}}) = \frac{1}{2} [\frac{\log(a^2-bc)}{n}+O(\frac{1}{n^2})]$ then $a^2-bc \neq 1$ div., $a^2-bc \neq 1 > 0$ con.
        \item $\sum (1-\frac{\log n}{n})^n$: $e^{n\log(1-\frac{\log n}{n})} = n^{-1}(1+o(1))$ div.
    \end{enumerate}
\end{prb*}

\begin{prb*}
    Following problems could be proved using Stirling's formula($n! = \sqrt{2\pi n}(\frac{n}{e})^n(1+O(\frac{1}{n}))$) or Sapagof test($\lim a_n = 0 \iff \sum 1-\frac{a_{n+1}}{a_n}$ div.): \\
    1. $\lim \frac{(2n)!}{4^n(n!)^2} = 0 \quad$ 
    2. $\lim \frac{n^n}{e^n n!} = 0 \quad$ 
    3. $\lim \frac{n!}{(x+1)\cdots(x+n)}=0\ (x>0)$
\end{prb*}

\begin{prb*}
    \checkmark 1. Given +-term $\sum a_n$ con., then it's possible that $\limsup na_n \geq 1$; 
    2.If $a_n$ also is monotonic, then $\limsup na_n = 0$.
\end{prb*}
\fbox{\begin{minipage}{0.95\textwidth}
    \begin{sol*}
        1. Core idea: previous section 1st problem solution 1.3). \\
        $a_n = \begin{cases} \frac{1}{2^n}\ n \neq k^2 \\ \frac{1}{n}\ n = k^2 \end{cases}$, $\limsup na_n = 1,\ \sum a_n < \infty$; 
        $a_n = \begin{cases} \frac{1}{2^n}\ n \neq k^3 \\ \frac{1}{\sqrt{n}}\ n = k^3 \end{cases}$, $\limsup na_n = \infty,\ \sum a_n < \infty$ (compared with $\sum \frac{1}{2^n}+\sum \frac{1}{n^p}\ (p>1)$); \\
        2. Using Cauchy's condensation Theorem/Cauchy's convergence Theorem.
    \end{sol*}
\end{minipage}}

\begin{prb*}
    Given +-term $\sum a_n$ con., $p>1$, then $\sum a_n^p$ con. If it's not +-term. does it still hold?
\end{prb*}
\fbox{\begin{minipage}{0.95\textwidth}
    \begin{sol*}
        A good example to illustrate Comparison Test is only applicable to +-term series. CounterExample: $\sum (-1)^n \frac{1}{\sqrt{n}}$. 
    \end{sol*}
\end{minipage}}

\begin{prb*}
    \checkmark Given $\{a_n\} \downarrow \to 0$, $b_n = a_n-2a_{n+1}+a_{n+2} \geq 0$, then $\sum nb_n = a_1$.
\end{prb*}
\fbox{\begin{minipage}{0.95\textwidth}
    \begin{sol*}
        Observing that $b_n = (a_n-a_{n+1})-(a_{n+1}-a_{n+2})$ and $\sum nb_n = \sum_{n=1}^{\infty} \sum_{k=n}^{\infty} b_k = a_1-\lim a_{n+1}-\lim n(a_{n+1}-a_{n+2})$, then i.e. problem is to prove $a_{n+1}-a_{n+2} = o(\frac{1}{n}),\ n\to\infty$ ($\Leftarrow$ $(a_n-a_{n+1}) \downarrow$ and $\sum (a_n-a_{n+1})$ con.) (Application of precious \checkmark problem)
    \end{sol*}
\end{minipage}}

\begin{prb*}
    If $\forall n,\ 0<a_n<a_{2n}+a_{2n+1}$, then $\sum a_n$ div.
\end{prb*}
\fbox{\begin{minipage}{0.95\textwidth}
    \begin{sol*}
        $0<a_n<a_{2n}+a_{2n+1}<a_{4n}+a_{4n+1}+a_{4n+2}+a_{4n+3}<\cdots<a_{2^mn}+\cdots+a_{2^m(n+1)-1}$ against to Cauchy's convergence Theorem.
    \end{sol*}
\end{minipage}}

\begin{prb*}
    Given $x>0$, then $\sum \frac{x^n}{(1+x)\cdots(1+x^n)}$ con.
\end{prb*}
\fbox{\begin{minipage}{0.95\textwidth}
    \begin{sol*}
        Easy: $S=\sum [\frac{1}{(1+x)\cdots(1+x^{n-1})}-\frac{1}{(1+x)\cdots(1+x^{n})}]=\frac{1}{1+x}+o(1)$
    \end{sol*}
\end{minipage}}

\begin{prb*}
    \checkmark Given $x>0$, test convergence of $\sum (2-x)(2-x^{\frac{1}{2}})\cdots(2-x^{\frac{1}{n}})$.
\end{prb*}
\fbox{\begin{minipage}{0.95\textwidth}
    \begin{sol*}
        1. $0<x\leq 1$, $(2-x)(2-x^{\frac{1}{2}})\cdots(2-x^{\frac{1}{n}}) \geq 1$ div. \\
        2. $x>1$, $\log \frac{a_{n}}{a_{n-1}} = \log (2-x^{\frac{1}{n}}) = \log (2-e^{\frac{\log x}{n}}) = -\frac{\log x}{n}+O(\frac{1}{n^2})$ $\Sigma: \log a_n = -\log x\log n+O(1)$ then $a_n = (\frac{1}{n})^{\log x}\cdot O(1)$: 1) $x>e$ con. 2) $x \leq e$ div.
    \end{sol*}
\end{minipage}}

\begin{prb*}
    Given +-term $\sum a_n$, $k\in\mathbb{Z}^{+}$, if $\lim \frac{a_{n+k}}{a_n} = l$, then $l<1$, $\sum a_n$ con.; $l>1$, $\sum a_n$ div. (Developement of d'Alembert Test)
\end{prb*}
\fbox{\begin{minipage}{0.95\textwidth}
    \begin{sol*}
        By definition, $\forall \varepsilon>0, \exists N_{\varepsilon}>0$ s.t. $\forall n \geq N_{\varepsilon}:\ l-\varepsilon<\frac{a_{n+k}}{a_n}<l+\varepsilon,\ l-\varepsilon<\frac{a_{n+(k+1)}}{a_{n+1}}<l+\varepsilon,\cdots,\ l-\varepsilon<\frac{a_{n+(2k-1)}}{a_{n+k-1}}<l+\varepsilon$, then 
        \begin{align*}
            \sum a_n &= \sum_{n=1}^{N} a_n + \sum_{i=0}^{k-1}\sum_{j=1}^{\infty} a_{n+jk+i} \\
            &> \sum_{n=1}^{N} a_n + \sum_{i=0}^{k-1}\sum_{j=1}^{\infty} (l-\varepsilon)^ja_{n+i} \\
            &< \sum_{n=1}^{N} a_n + \sum_{i=0}^{k-1}\sum_{j=1}^{\infty} (l+\varepsilon)^ja_{n+i}
        \end{align*}
    \end{sol*}
\end{minipage}}

\begin{prb*}
    Given $0<a_1<\frac{\pi}{2}$, $a_n=\sin a_{n-1}$, then test the convergence of $\sum a_n^p$.
\end{prb*}
\fbox{\begin{minipage}{0.95\textwidth}
    \begin{sol*}
        \fbox{\begin{minipage}{0.78\textwidth}
            Lemma. Let $x_{0}\in(0,\pi)$, $x_n=\sin x_{n-1}$, then $x_{n}=\frac{3^{1/2}}{n^{1/2}}-\frac{3^{3/2}\ln n}{10n^{3/2}}+\mathcal O(n^{-3/2})$.
        \end{minipage}}
        Then $p>2$, $\sum a_n^p$ con.; $p<2$, $\sum a_n^p$ div.
    \end{sol*}
\end{minipage}}
\newline\newline
\begin{nt*}
    Readers could try to prove the Lemma.
\end{nt*}

\subsubsection{General term Series}
Strategies of study general term series: \\
1. we could first study +-term $\sum |a_n|$; \\
2. from simple to complex, try various of Tests: Leibniz, Dirichlet, Abel, so on; \\
3. the equivalence test (order estimation) is usually not valid for series with changing signs (of course it still holds when there is abs. con.), but we can still use it to roughly analyze the asymptotic behavior of the general term of the series. With the help of the rearrangement properties of the general term series that meet the conditions, such as Riemann's rearrangement theorem, we can still partially use the method of estimating order; \\
4. if the above methods do not work, please make use of various trigonometric function conversion formulas, when dealing with summation problems involving trigonometric functions. \\
Here we first introduce some theorems and tips/tools, then dive into problem solving: \\
\fbox{\begin{minipage}{0.95\textwidth}
    \begin{thm}
        (Riemann Seires Theorem) \\
        Given $\{a_n\}$, and $\sum_{n=1}^{\infty} a_n$ is con. con., $M>0$, then there $\exists$ a permutation $\sigma$ s.t. $\sum_{n=1}^{\infty} a_{\sigma(n)} = M$.
        There also $\exists$ a permutation $\sigma$ s.t. $\sum_{n=1}^{\infty} a_{\sigma(n)} = \infty$.
        The sum can also be rearranged to diverge to $-\infty$ or to fail to approach any limit, finite or infinite.
    \end{thm}
\end{minipage}}
\newline\newline
